\section{Is High Compactness a Legal Outlier?}
\label{sec:outlier}

\subsection{The Two-Sided Risk}

Courts have typically treated low compactness as a marker of gerrymandering:
a district that ``snakes'' through a city to collect partisan voters is suspect.
This is the logic of \shaw{} and its progeny, which applied strict scrutiny to
``extremely irregular'' racial gerrymanders \citep{scotus1993shaw}.

But there is a less-discussed risk on the other side: a plan that is
\emph{unusually compact} may have achieved that compactness by packing
minority or partisan voters into a small number of districts.  In sorted
states, a maximally compact plan tends to pack Democratic voters into
urban districts, producing an efficient Republican majority.  A plan at the
99th percentile of compactness in a sorted state would warrant scrutiny on
this basis.

\begin{keybox}
\textbf{The two-sided compactness test:} A legally defensible plan should
be more compact than average (to demonstrate non-gerrymandering by the
traditional test) but not at the extreme right tail of the compactness
distribution (to avoid the packing accusation in sorted states).
The $61$st--$72$nd percentile range satisfies both conditions.
\end{keybox}

\subsection{Quantifying the Outlier Threshold}

Using the distribution parameterisation from Section~\ref{sec:distribution},
the 95th percentile of mean Polsby-Popper score is approximately
$0.42$--$0.46$ for the states studied.  The \ar{} plan's values
($0.362$--$0.412$) are below this threshold in all states.

More precisely:

\begin{center}
\begin{tabular}{lccc}
\toprule
State & $\bar\pp_{\ar}$ & Ensemble 95th pct & Outlier? \\
\midrule
NC & 0.412 & 0.443 & No \\
WI & 0.388 & 0.427 & No \\
GA & 0.395 & 0.436 & No \\
PA & 0.371 & 0.412 & No \\
\bottomrule
\end{tabular}
\end{center}

No \ar{} state plan exceeds the 95th percentile compactness threshold.
The plan is compact but not an outlier.

\subsection{The Partisan-Compact Correlation in Ensemble Plans}

Within the \recom{} ensemble, plans at the 95th percentile of compactness
tend to produce more Republican seats than plans at the 50th percentile of
compactness, in sorted states.  This is consistent with the Rodden effect:
compact plans in sorted states pack Democratic voters.

The \ar{} plan, at the 65th--72nd percentile, experiences a mild version
of this effect.  It is more compact than average, and in Georgia it produces
one fewer Democratic seat than the ensemble median.  But it is not at the
extreme where the packing accusation becomes compelling.
